\section{Introduction}
\label{sec:introduction}

People use multiple large language models (LLMs), but each interface typically
receives only a local account setting, a manually repeated instruction, or a
provider-specific memory. The resulting problem is not only remembering what a
person said. It is representing \emph{how that person wants a response to be
constructed for the task at hand} and transferring that representation without
assuming that expertise or style in one domain applies everywhere.

We call this proposed layer \systemname{}. Its design goals are portability across
model providers, contextual selection, and transparent user control. A response
preference is an instruction-like behavioral constraint---such as desired depth,
format, or use of examples---rather than a demographic label or an assertion
about identity. The current request always outranks a stored default.

The primary research question is whether a dynamically learned,
domain-conditioned, provider-independent preference representation improves
subjective response quality across heterogeneous model families. We further ask
whether it improves on static profiles and history retrieval, whether learning is
reliable under drift and adversarial evidence, and whether structured profiles
reduce context overhead.

The contributions of the present artifact are:
\begin{enumerate}[leftmargin=*]
  \item a formal, inspectable \emph{preference policy} representation and
  applicability problem that separates response behavior from conversation memory;
  \item a behavioral-conformance framing of portability: an unchanged profile
  must induce the intended scoped behavior across model families, not merely
  serialize successfully;
  \item a benchmark and human-study protocol targeting domain isolation,
  preference learning, context efficiency, drift, and incorrect inference; and
  \item a local-first reference implementation and reproducible comparison of five
  personalization conditions, including 128 real local-model generations and an
  observed prompt-compression/output-expansion trade-off.
\end{enumerate}

The representation and pipeline are implemented contributions; the pilot is
exploratory systems evidence. It does not establish the primary subjective-quality
hypothesis. Irrelevant, stale, or incorrect preferences can still reduce quality
or undermine user control.
